% !TEX root = ../lectures.tex

{\large\bfseries Calculus 2}

\section{Lecture 7}
\subsection*{Power Series}

\begin{definition}\label{def:7.1}
A functional series of the form
\[
\sum_{k=0}^{\infty}a_k(x-x_0)^k
=a_0+a_1(x-x_0)+a_2(x-x_0)^2+\cdots+a_k(x-x_0)^k+\cdots
\tag{1}\label{eq:7:1:1}
\]
is called a \emph{power series centered at $x_0$}.

$x_0\in\R$ is called the \emph{center} of the series.

$a_k\in\R$, $k\in\N\cup\{0\}$ are the \emph{coefficients} of the series.

Since by the linear substitution $t=x-x_0$, the power series (1) can be rewritten as a power series centered at $0$:
\[
\sum_{k=0}^{\infty}a_kt^k
=a_0+a_1t+a_2t^2+\cdots+a_kt^k+\cdots,
\tag{2}\label{eq:7:2:1}
\]
without loss of generality, we can study only power series of the form (2).
\end{definition}

\begin{definition}\label{def:7.2}
Let
\[
D=\left\{x\in\R\ \middle|\ \sum_{k=0}^{\infty}a_kx^k\text{ converges}\right\}\subseteq\R.
\]
That is, $D$ is the \emph{set of all points of convergence} of $\sum_{k=0}^{\infty}a_kx^k$.

Then,
\[
R=\sup\{|x|\mid x\in D\}
\]
is called the \emph{radius of convergence} of the power series $\sum_{k=0}^{\infty}a_kx^k$.
The set $D$ is called the \emph{interval of convergence} of the power series $\sum_{k=0}^{\infty}a_kx^k$.

Note that, since $0\in D$, we have $D\neq\varnothing$, which implies that
$R=\sup\{|x|\mid x\in D\}$ exists.
\end{definition}

\begin{claim}\label{claim:7.1}
Let $\sum_{k=0}^{\infty}a_kx^k$ be a power series and let $R$ be its radius of convergence.
Then:
\begin{enumerate}
    \item\label{item:claim:7.1:1} $\forall x\in\R:\ |x|<R$, the series $\sum_{k=0}^{\infty}a_kx^k$ converges \emph{absolutely};
    \item\label{item:claim:7.1:2} $\forall x\in\R:\ |x|>R$, the series $\sum_{k=0}^{\infty}a_kx^k$ \emph{diverges}.
\end{enumerate}
\end{claim}

\begin{proof}
1) Due to \hyperref[def:3.1]{Def 3.1}, it suffices to show that the series
\[
\sum_{k=0}^{\infty}|a_kx^k|
\]
converges.

Since
\[
|x|<R=\sup\{|y|\mid y\in D\},
\qquad
D=\left\{y\in\R\ \middle|\ \sum_{k=0}^{\infty}a_ky^k\text{ converges}\right\},
\]
by the definition of the supremum, there exists $y\in D$ such that $|x|<|y|$.
Since $y\in D$, the series $\sum_{k=0}^{\infty}a_ky^k$ converges, therefore, due to \hyperref[claim:1.5]{Claim 1.5}, we have
\[
\lim_{k\to\infty}a_ky^k=0,
\]
which implies that
\[
\exists M>0\ \forall k\in\N\cup\{0\}\to |a_ky^k|<M.
\]
Therefore, for every $k\in\N\cup\{0\}$ we have
\[
|a_kx^k|
=|a_ky^k|\left|\frac{x}{y}\right|^k
<M\left|\frac{x}{y}\right|^k
=Mq^k,
\]
where
\[
0\leq |x|<|y|\quad\Longrightarrow\quad q=\left|\frac{x}{y}\right|\in[0,1).
\]
The latter implies that the series
\[
\sum_{k=0}^{\infty}Mq^k
\]
converges. Therefore, due to \hyperref[claim:5.7]{Claim 5.7} [The Weierstrass M-test], the series
\[
\sum_{k=0}^{\infty}|a_kx^k|
\]
converges.

2) Assume that $|x|>R$ and $\sum_{k=0}^{\infty}a_kx^k$ converges.
The latter implies (see \hyperref[def:7.2]{Def 7.2}) that $x\in D$.
Therefore, since
\[
R=\sup\{|y|\mid y\in D\},
\]
we have $|x|\leq R$, which contradicts $|x|>R$.
Thus, our assumption is false and the series $\sum_{k=0}^{\infty}a_kx^k$ diverges.
\end{proof}

\begin{claim}[The Cauchy--Hadamard Formula]\label{claim:7.2}
For any power series $\sum_{k=0}^{\infty}a_kx^k$, its radius of convergence $R$ is given by
\[
\frac{1}{R}
=
\overline{\lim}_{k\to\infty}|a_k|^{1/k}
=
\lim_{k\to\infty}\sup_{n\geq k}|a_n|^{1/n}.
\tag{*}\label{eq:7:star:1}
\]
\end{claim}

\begin{proof}
Recall that, by \hyperref[def:7.2]{Def 7.2}, we have
\[
D=\left\{x\in\R\ \middle|\ \sum_{k=0}^{\infty}a_kx^k\text{ converges}\right\}
\qquad\text{and}\qquad
R=\sup\{|x|\mid x\in D\}.
\]
It is clear that $0\in D$.
Let $x\neq0$, then we have
\[
L
=
\overline{\lim}_{k\to\infty}|a_kx^k|^{1/k}
=
\overline{\lim}_{k\to\infty}|a_k|^{1/k}|x|
=
|x|\tau,
\tag{1}\label{eq:7:1:2}
\]
where
\[
\tau=\overline{\lim}_{k\to\infty}|a_k|^{1/k}.
\]
\hyperref[eq:7:1:2]{Equality (1)} holds true since $x$ does not depend on $k$ and $|x|>0$ (see \hyperref[claim:7.4]{Claim 7.4}).

Then, due to \hyperref[claim:2.9]{Claim 2.9} [The Root Test], the series $\sum_{k=0}^{\infty}|a_kx^k|$ converges if $L=|x|\tau<1$ and diverges if $L=|x|\tau>1$.

\begin{itemize}
    \item If $\tau=0$, then $|x|\tau=0<1$, $\forall x\in\R$, thus $D=\R$ and $R=+\infty$; that is, \hyperref[eq:7:star:1]{Formula (*)} holds true if $\tau=0$.

    \item If $\tau=+\infty$, then $|x|\tau=+\infty>1$, $\forall x\in\R\setminus\{0\}$, thus $D=\{0\}$ and $R=0$; that is, \hyperref[eq:7:star:1]{Formula (*)} holds true if $\tau=+\infty$.

    \item If $\tau\in(0,+\infty)$, then
    \[
    |x|\tau<1
    \Longleftrightarrow
    |x|<\frac{1}{\tau},
    \]
    hence
    \[
    \left\{x\in\R\ \middle|\ |x|<\frac{1}{\tau}\right\}\subseteq D,
    \]
    which implies that $R\geq\frac{1}{\tau}$.

    Let us assume that $R>\frac{1}{\tau}$. Then, due to \hyperref[def:7.2]{Def 7.2}, there is $x\in D$ such that
    \[
    R>|x|>\frac{1}{\tau}.
    \]
    Inequality $|x|>\frac{1}{\tau}$ implies that the series $\sum_{k=0}^{\infty}|a_kx^k|$ diverges (since $L=|x|\tau>1$), and inequality $R>|x|$ implies that the series $\sum_{k=0}^{\infty}a_kx^k$ converges (due to \hyperref[claim:7.1]{Claim 7.1}).
    Thus, the assumption is false.
\end{itemize}
Therefore, $R=\frac{1}{\tau}$ and \hyperref[eq:7:star:1]{Formula (*)} holds true if $\tau\in(0,+\infty)$.
\end{proof}

\textbf{Note} that \hyperref[claim:7.1]{Claim 7.1} tells nothing about the behavior of the series $\sum_{k=0}^{\infty}a_kx^k$ at the boundary points $x=R$ and $x=-R$, if $R\in(0,+\infty)$.

In fact, all logically admissible possibilities can occur:
\begin{itemize}
    \item the series can be absolutely convergent at $x=R$ and $x=-R$;

    for example:
    \[
    \sum_{k=1}^{\infty}\frac{1}{k^2}x^k,
    \qquad R=1,
    \qquad D=[-1,1];
    \]

    \item the series can be divergent at $x=R$ and $x=-R$;

    for example:
    \[
    \sum_{k=1}^{\infty}\frac{1}{k}x^{2k},
    \qquad R=1,
    \qquad D=(-1,1);
    \]

    \item the series can be divergent at $x=R$ and conditionally convergent at $x=-R$;

    for example:
    \[
    \sum_{k=1}^{\infty}\frac{1}{k}x^k,
    \qquad R=1,
    \qquad D=[-1,1);
    \]

    \item the series can be conditionally convergent at $x=R$ and divergent at $x=-R$;

    for example:
    \[
    \sum_{k=1}^{\infty}\frac{(-1)^k}{k}x^k,
    \qquad R=1,
    \qquad D=(-1,1].
    \]
\end{itemize}

\begin{claim}[On Uniform Convergence of Power Series]\label{claim:7.3}
Let $\sum_{k=0}^{\infty}a_kx^k$ be a power series with the radius of convergence $R>0$.
Then:
\begin{enumerate}
    \item\label{item:claim:7.3:1} the series converges \emph{uniformly} on every interval
    \[
    I=[a,b]\subset(-R,R);
    \]

    \item\label{item:claim:7.3:2} if the series \emph{converges at $x=R$}, then, it converges \emph{uniformly} on every interval
    \[
    I=[a,R]\subset(-R,R];
    \]

    \item\label{item:claim:7.3:3} if the series \emph{diverges at $x=R$}, then, it converges \emph{non-uniformly} on every interval
    \[
    I=[a,R)\subset(-R,R).
    \]
\end{enumerate}
\end{claim}

\begin{proof}
1) Let
\[
m=\max\{|a|,|b|\}.
\]
It is clear that $0\leq m<R$.
Note that:
\begin{itemize}
    \item $\forall k\in\N\cup\{0\}$, $\forall x\in I=[a,b]$,
    \[
    |a_kx^k|\leq |a_k|m^k=:M_k;
    \]

    \item since $|m|<R$, due to \hyperref[item:claim:7.1:1]{Item 1)} of \hyperref[claim:7.1]{Claim 7.1}, the series
    \[
    \sum_{k=0}^{\infty}M_k
    =
    \sum_{k=0}^{\infty}|a_k|m^k
    \]
    converges.
\end{itemize}
Thus, due to \hyperref[claim:5.7]{Claim 5.7} [The Weierstrass M-test], the series $\sum_{k=0}^{\infty}a_kx^k$ converges uniformly on $[a,b]\subset(-R,R)$.

2) Due to \hyperref[item:claim:7.3:1]{Item 1)} and \hyperref[claim:6.6]{Claim 6.6}, it suffices to prove the claim for $I=[0,R]$.

Let
\[
f_k(x)=a_kR^k
\qquad\text{and}\qquad
g_k(x)=\left(\frac{x}{R}\right)^k,
\qquad x\in[0,R],\quad \forall k\in\N\cup\{0\}.
\]
Note that:
\begin{itemize}
    \item the series
    \[
    \sum_{k=0}^{\infty}f_k(x)
    =
    \sum_{k=0}^{\infty}a_kR^k
    \]
    converges uniformly on $[0,R]$.

    Indeed, the series converges since this is one of the conditions of the claim; since the series $\sum_{k=0}^{\infty}a_kR^k$ does not effectively depend on $x$, we can say, due to \hyperref[def:5.3]{Def 5.3}, that it converges uniformly on $[0,R]$.

    \item $\forall k\in\N\cup\{0\}$, $\forall x\in[0,R]$,
    \[
    \left(\frac{x}{R}\right)^k\geq\left(\frac{x}{R}\right)^{k+1},
    \]
    i.e. $g_k(x)\geq g_{k+1}(x)$;

    \item $\forall k\in\N\cup\{0\}$, $\forall x\in[0,R)$,
    \[
    \left|\frac{x}{R}\right|<1,
    \]
    i.e. $g_k(x)\downarrow0$.
\end{itemize}
Thus, due to \hyperref[claim:5.9]{Claim 5.9} [The Abel Test for the UC], the series
\[
\sum_{k=0}^{\infty}f_k(x)g_k(x)
=
\sum_{k=0}^{\infty}a_kR^k\left(\frac{x}{R}\right)^k
=
\sum_{k=0}^{\infty}a_kx^k
\]
converges uniformly on $[0,R]$.

3) The pointwise convergence of $\sum_{k=0}^{\infty}a_kx^k$ on $[a,R)$ follows from \hyperref[item:claim:7.1:1]{Item 1)} of \hyperref[claim:7.1]{Claim 7.1}.
Let us assume that the series
\[
\sum_{k=0}^{\infty}a_kx^k
=
\lim_{n\to\infty}f_n(x),
\qquad
f_n(x)=\sum_{k=0}^{n}a_kx^k,
\]
converges uniformly on $[a,R)$.
Note that:
\begin{itemize}
    \item $\lambda=R$ is a finite limit point of the set $[a,R)$;
    \item $\forall n\in\N$, we have
    \[
    \lim_{x\to R}f_n(x)=\sum_{k=0}^{n}a_kR^k\in\R
    \]
    (since $n$ is finite).
\end{itemize}
Thus, due to \hyperref[item:claim:6.1:1]{Item 1)} of \hyperref[claim:6.1]{Claim 6.1} [On Changing the Order of Two Limits], the limit
\[
\lim_{n\to\infty}\sum_{k=0}^{n}a_kR^k
=
\sum_{k=0}^{\infty}a_kR^k
\]
exists. Since this contradicts the conditions of the claim, our assumption is false and the series $\sum_{k=0}^{\infty}a_kx^k$ converges non-uniformly on $[a,R)\subset(-R,R)$.
\end{proof}

\begin{claim}\label{claim:7.4}
Let $\{a_k\}_{k=1}^{\infty}$ and $\{b_k\}_{k=1}^{\infty}$ be two numerical sequences.
If
\[
\exists\lim_{k\to\infty}a_k=\alpha>0,
\]
then
\[
\overline{\lim}_{k\to\infty}(a_kb_k)
=
\alpha\cdot\overline{\lim}_{k\to\infty}b_k.
\]
\end{claim}

\begin{proof}
See Calculus-1.
\end{proof}

\begin{claim}[On Derivative of Power Series]\label{claim:7.5}
Let
\[
S(x)=\sum_{k=0}^{\infty}a_kx^k
\]
on $(-R,R)$, where $R>0$ is the radius of convergence of the power series.
Then:
\begin{enumerate}
    \item\label{item:claim:7.5:1} the radius of convergence of the series
    \[
    \sum_{k=1}^{\infty}ka_kx^{k-1}
    \]
    is equal to $R$;

    \item\label{item:claim:7.5:2} for every $x\in(-R,R)$, we have
    \[
    S'(x)
    =
    \left(\sum_{k=0}^{\infty}a_kx^k\right)'
    =
    \sum_{k=1}^{\infty}ka_kx^{k-1}.
    \tag{3}\label{eq:7:3:1}
    \]
\end{enumerate}
\end{claim}

\begin{proof}
1) Note that:
\begin{itemize}
    \item the series
    \[
    \sum_{k=1}^{\infty}ka_kx^{k-1}
    \qquad\text{and}\qquad
    \sum_{k=1}^{\infty}ka_kx^k
    =
    x\sum_{k=1}^{\infty}ka_kx^{k-1}
    \]
    have the same interval of convergence, therefore, due to \hyperref[def:7.2]{Def 7.2}, they have the same radius of convergence;

    \item since
    \[
    \lim_{k\to\infty}k^{1/k}=1>0,
    \]
    due to \hyperref[claim:7.4]{Claim 7.4}, we have the equality
    \[
    \begin{aligned}
    \overline{\lim}_{k\to\infty}|ka_k|^{1/k}
    &=\overline{\lim}_{k\to\infty}\left(k^{1/k}|a_k|^{1/k}\right)\\
    &=1\cdot\overline{\lim}_{k\to\infty}|a_k|^{1/k}
    =\overline{\lim}_{k\to\infty}|a_k|^{1/k}.
    \end{aligned}
    \]
    which, due to \hyperref[claim:7.2]{Claim 7.2} [The Cauchy--Hadamard Formula], implies that the series $\sum_{k=1}^{\infty}ka_kx^k$ and $\sum_{k=0}^{\infty}a_kx^k$ have the same radius of convergence.
\end{itemize}
Therefore, the series $\sum_{k=1}^{\infty}ka_kx^{k-1}$ and $\sum_{k=0}^{\infty}a_kx^k$ have the same radius of convergence and \hyperref[item:claim:7.5:1]{Item 1)} is proved.

2) Let $x_0\in(-R,R)$ be arbitrary. Then it ought to be clear that there is an open interval $J=(-r,r)$ such that $x_0\in J$ and $0<r<R$.

Let
\[
u_k(x)=a_kx^k,\qquad \forall k\in\N\cup\{0\},
\]
and
\[
f_n(x)=\sum_{k=0}^{n}u_k(x),\qquad \forall n\in\N\cup\{0\}.
\]
Then note that:
\begin{enumerate}[label=\alph*)]
    \item $\forall k\in\N\cup\{0\}$, $u_k(x)=a_kx^k$ is differentiable on $J=(-r,r)$;

    \item the series $\sum_{k=0}^{\infty}u_k(x)$ converges at $x=0\in J$;

    \item the series
    \[
    \sum_{k=0}^{\infty}u_k'(x)
    =
    \sum_{k=1}^{\infty}ka_kx^{k-1}
    \]
    converges uniformly on $J$.
\end{enumerate}
Indeed, due to \hyperref[item:claim:7.5:1]{Item 1)}, the radius of convergence of the series $\sum_{k=1}^{\infty}ka_kx^{k-1}$ is equal to $R>0$.
Therefore, due to \hyperref[item:claim:7.3:1]{Item 1)} of \hyperref[claim:7.3]{Claim 7.3}, it converges uniformly on the interval $[-r,r]\subset(-R,R)$, which implies the uniform convergence of the series on $(-r,r)$.

That is, all conditions of \hyperref[claim:6.13]{Claim 6.13} [Term-by-term Diff. of Series] are satisfied and \hyperref[eq:7:3:1]{Equality (3)} holds true $\forall x\in(-r,r)$ due to \hyperref[item:claim:6.13:2]{Item 2)} of \hyperref[claim:6.13]{Claim 6.13}.
Then, since $x_0\in(-R,R)$ is arbitrary, \hyperref[eq:7:3:1]{Equality (3)} holds true $\forall x\in(-R,R)$ and \hyperref[item:claim:7.5:2]{Item 2)} is proved.
\end{proof}

\textbf{Note} that the series $\sum_{k=0}^{\infty}a_kx^k$ and $\sum_{k=1}^{\infty}ka_kx^{k-1}$ have the same radius of convergence, but the series may have different intervals of convergence (see \hyperref[def:7.2]{Def 7.2}).

For example, if $a_0=0$, $a_k=\frac{1}{k^2}$, $\forall k\in\N$, we have
\[
\sum_{k=0}^{\infty}a_kx^k
=
\sum_{k=1}^{\infty}\frac{1}{k^2}x^k
\qquad\text{with }R=1\text{ and }D=[-1,1],
\]
and
\[
\sum_{k=1}^{\infty}ka_kx^{k-1}
=
\sum_{k=1}^{\infty}\frac{1}{k}x^{k-1}
\qquad\text{with }R=1\text{ and }D=[-1,1).
\]

\begin{claim}[Analyticity of Power Series]\label{claim:7.6}
Let
\[
S(x)=\sum_{k=0}^{\infty}a_kx^k
\]
on $(-R,R)$, where $R>0$ is the radius of convergence of the series.
Then $S\in C^{\infty}(-R,R)$ and, for $\forall n\in\N$, $\forall x\in(-R,R)$, we have
\[
S^{(n)}(x)
=
\sum_{k=n}^{\infty}k(k-1)\cdots(k-n+1)a_kx^{k-n},
\]
where $C^{\infty}(-R,R)$ is the set of all infinitely differentiable functions on $(-R,R)$ and $S^{(n)}(x)$ is the $n$-th derivative of $S(x)$.
\end{claim}

\begin{proof}
Follows directly from \hyperref[claim:7.5]{Claim 7.5} via mathematical induction on $n\in\N$.
\end{proof}

\begin{claim}[On Antiderivative of Power Series]\label{claim:7.7}
Let
\[
S(x)=\sum_{k=0}^{\infty}a_kx^k
\]
on $(-R,R)$, where $R>0$ is the radius of convergence of the series.
Then:
\begin{enumerate}
    \item\label{item:claim:7.7:1} the radius of convergence of the series
    \[
    \sum_{k=0}^{\infty}\frac{a_k}{k+1}x^{k+1}
    \]
    is equal to $R$;

    \item\label{item:claim:7.7:2} for every $x\in(-R,R)$, we have
    \[
    \int_0^x S(t)\,dt
    =
    \int_0^x\left(\sum_{k=0}^{\infty}a_kt^k\right)dt
    =
    \sum_{k=0}^{\infty}\frac{a_k}{k+1}x^{k+1}
    =
    \sum_{k=0}^{\infty}\left(\int_0^x a_kt^k\,dt\right).
    \tag{4}\label{eq:7:4:1}
    \]
\end{enumerate}
\end{claim}

\begin{proof}
1) Note that:
\begin{itemize}
    \item the series
    \[
    \sum_{k=0}^{\infty}\frac{a_k}{k+1}x^{k+1}
    =
    x\sum_{k=0}^{\infty}\frac{a_k}{k+1}x^k
    \]
    and the series
    \[
    \sum_{k=0}^{\infty}\frac{a_k}{k+1}x^k
    \]
    have the same interval of convergence, therefore, due to \hyperref[def:7.2]{Def 7.2}, they have the same radius of convergence;

    \item since
    \[
    \lim_{k\to\infty}\left(\frac{1}{k+1}\right)^{1/k}=1>0,
    \]
    due to \hyperref[claim:7.4]{Claim 7.4}, we have the equality
    \[
    \begin{aligned}
    \overline{\lim}_{k\to\infty}\left|\frac{1}{k+1}a_k\right|^{1/k}
    &=\overline{\lim}_{k\to\infty}\left(\left(\frac{1}{k+1}\right)^{1/k}|a_k|^{1/k}\right)\\
    &=1\cdot\overline{\lim}_{k\to\infty}|a_k|^{1/k}
    =\overline{\lim}_{k\to\infty}|a_k|^{1/k}.
    \end{aligned}
    \]
    which, due to \hyperref[claim:7.2]{Claim 7.2} [The Cauchy--Hadamard Formula], implies that the series $\sum_{k=0}^{\infty}\frac{a_k}{k+1}x^k$ and $\sum_{k=0}^{\infty}a_kx^k$ have the same radius of convergence.
\end{itemize}
Therefore, the series $\sum_{k=0}^{\infty}\frac{a_k}{k+1}x^{k+1}$ and $\sum_{k=0}^{\infty}a_kx^k$ have the same radius of convergence and \hyperref[item:claim:7.7:1]{Item 1)} is proved.

2) Let $x\in(-R,R)$, $I=[0,x]\subset(-R,R)$,
\[
u_k(t)=a_kt^k,
\qquad \forall k\in\N\cup\{0\},
\]
and
\[
f_n(t)=\sum_{k=0}^{n}u_k(t).
\]
Then note that:
\begin{enumerate}[label=\alph*)]
    \item $\forall k\in\N\cup\{0\}$, $u_k(t)=a_kt^k$ is Riemann integrable on $I=[0,x]$;

    \item the series $\sum_{k=0}^{\infty}a_kt^k$ converges uniformly on $I=[0,x]$.
\end{enumerate}
The latter holds since $I=[0,x]\subset(-R,R)$ and due to \hyperref[item:claim:7.3:1]{Item 1)} of \hyperref[claim:7.3]{Claim 7.3}.
Thus, all conditions of \hyperref[claim:6.12]{Claim 6.12} [Term-by-term Int. of Series] are satisfied and, for every $x\in(-R,R)$, we have
\[
\int_0^x S(t)\,dt
=
\int_0^x\left(\sum_{k=0}^{\infty}a_kt^k\right)dt
\overset{1}{=}
\sum_{k=0}^{\infty}\left(\int_0^x a_kt^k\,dt\right)
\overset{2}{=}
\sum_{k=0}^{\infty}\frac{a_k}{k+1}x^{k+1}.
\]
Here, 1) is due to \hyperref[claim:6.12]{Claim 6.12} and 2) follows from
\[
\int_0^x a_kt^k\,dt
=
\left.\frac{a_k}{k+1}t^{k+1}\right|_0^x
=
\frac{a_k}{k+1}x^{k+1}.
\]
Thus, \hyperref[eq:7:4:1]{Equality (4)} is proved.
\end{proof}

\textbf{Note} that \hyperref[item:claim:7.7:1]{Item 1)} of \hyperref[claim:7.7]{Claim 7.7} tells us that the series $\sum_{k=0}^{\infty}a_kx^k$ and $\sum_{k=0}^{\infty}\frac{a_k}{k+1}x^{k+1}$ have the same radius of convergence, but the series may have different intervals of convergence (see \hyperref[def:7.2]{Def 7.2}).

For example, if $a_0=0$, $a_k=\frac{1}{k}$, $\forall k\in\N$, we have
\[
\sum_{k=0}^{\infty}a_kx^k
=
\sum_{k=1}^{\infty}\frac{1}{k}x^k
\qquad\text{with }R=1\text{ and }D=[-1,1),
\]
and
\[
\sum_{k=0}^{\infty}\frac{a_k}{k+1}x^{k+1}
=
\sum_{k=1}^{\infty}\frac{1}{k(k+1)}x^{k+1}
\qquad\text{with }R=1\text{ and }D=[-1,1].
\]

\begin{claim}[On Equality of Power Series]\label{claim:7.8}
Let
\[
f(x)=\sum_{k=0}^{\infty}a_kx^k\quad\text{on }(-R_1,R_1)
\]
and
\[
g(x)=\sum_{k=0}^{\infty}b_kx^k\quad\text{on }(-R_2,R_2),
\]
where $R_1>0$ and $R_2>0$ are radii of convergence of the power series $\sum_{k=0}^{\infty}a_kx^k$ and $\sum_{k=0}^{\infty}b_kx^k$ respectively.
If there is a numerical sequence $\{x_n\}_{n=1}^{\infty}$ such that:
\begin{enumerate}[label=\alph*)]
    \item\label{item:claim:7.8:a} $x_n\in(-r,r)\setminus\{0\}$, $\forall n\in\N$, where
    \[
    r=\min\{R_1,R_2\}>0;
    \]
    \item\label{item:claim:7.8:b} $\displaystyle\lim_{n\to\infty}x_n=0$;
    \item\label{item:claim:7.8:c} $f(x_n)=g(x_n)$ for every $n\in\N$,
\end{enumerate}
then $a_k=b_k$, $\forall k\in\N\cup\{0\}$.
\end{claim}

\begin{proof}
Let
\[
c_k=a_k-b_k,
\qquad \forall k\in\N\cup\{0\}.
\]
Then, for every $x\in(-r,r)$, we have
\[
h(x):=\sum_{k=0}^{\infty}c_kx^k
=
\sum_{k=0}^{\infty}(a_k-b_k)x^k
\overset{1}{=}
\sum_{k=0}^{\infty}a_kx^k-
\sum_{k=0}^{\infty}b_kx^k
=f(x)-g(x),
\tag{**}\label{eq:7:starstar:1}
\]
where \hyperref[eq:7:1:2]{Equality 1}) holds since both series converge due to \hyperref[item:claim:7.1:1]{Item 1)} of \hyperref[claim:7.1]{Claim 7.1}.

From \hyperref[eq:7:starstar:1]{Formula (**)} it follows that:
\begin{itemize}
    \item the series
    \[
    h(x)=\sum_{k=0}^{\infty}c_kx^k
    \]
    has the radius of convergence $R\geq r>0$;

    \item
    \[
    h(x_n)=0,
    \qquad \forall n\in\N;
    \tag{5}\label{eq:7:5:1}
    \]

    \item it suffices to prove that
    \[
    c_k:=a_k-b_k=0,
    \qquad \forall k\in\N\cup\{0\}.
    \]
\end{itemize}

Let us assume that there is $k\in\N\cup\{0\}$ such that $c_k\neq0$.
Then
\[
m=\min\{k\in\N\cup\{0\}\mid c_k\neq0\}
\]
is well-defined and, for every $x\in(-r,r)$, we have
\begin{align*}
h(x)
&=\sum_{k=0}^{\infty}c_kx^k
=\sum_{k=0}^{m-1}c_kx^k+\sum_{k=m}^{\infty}c_kx^k\\
&=0+\sum_{k=m}^{\infty}c_kx^k
=x^m\sum_{k=0}^{\infty}c_{k+m}x^k.
\end{align*}
Here, the first finite sum is $0$ since $c_k=0$ for every $k\in\{0,1,\ldots,m-1\}$.

It ought to be clear that the series
\[
\widehat h(x):=\sum_{k=0}^{\infty}c_{k+m}x^k
\]
and
\[
h(x)=x^m\sum_{k=0}^{\infty}c_{m+k}x^k=x^m\widehat h(x)
\]
have the same interval of convergence, therefore, due to \hyperref[def:7.2]{Def 7.2}, the series have the same radius of convergence $R\geq r>0$.

Since, due to \hyperref[item:claim:7.8:a]{Item a)}, $x_n\neq0$ for every $n\in\N$, the equalities
\[
h(x)=x^m\widehat h(x),
\qquad \forall x\in(-r,r),
\]
and \hyperref[eq:7:5:1]{Formula (5)} imply that
\[
\widehat h(x_n)=0,
\qquad \forall n\in\N.
\]
Now note that
\[
c_m
\overset{1}{=}
\sum_{k=0}^{\infty}c_{k+m}0^k
\overset{2}{=}
\widehat h(0)
\overset{3}{=}
\widehat h\left(\lim_{n\to\infty}x_n\right)
\overset{4}{=}
\lim_{n\to\infty}\widehat h(x_n)
\overset{5}{=}
\lim_{n\to\infty}0
=0.
\]
Indeed:
\begin{enumerate}
    \item since $0^0:=1$ and $0^k=0$, $\forall k\in\N$;
    \item by definition, $\widehat h(x)=\sum_{k=0}^{\infty}c_{m+k}x^k$, $\forall x\in(-r,r)$, $r>0$;
    \item $\lim_{n\to\infty}x_n=0$ due to \hyperref[item:claim:7.8:b]{Item b)};
    \item due to \hyperref[item:claim:7.5:2]{Item 2)} of \hyperref[claim:7.5]{Claim 7.5}, $\widehat h$ is continuous on $(-r,r)$; since $\lim_{n\to\infty}x_n=0\in(-r,r)$, this implies that
    \[
    \lim_{n\to\infty}\widehat h(x_n)
    =
    \widehat h\left(\lim_{n\to\infty}x_n\right);
    \]
    \item due to \hyperref[eq:7:5:1]{Equality (5)}, we have $\widehat h(x_n)=0$, $\forall n\in\N$.
\end{enumerate}
That is, we arrived at a contradiction (with $c_m\neq0$).
Therefore, the assumption is false and $c_k=0$, $\forall k\in\N\cup\{0\}$.
\end{proof}

Now we can easily prove \hyperref[claim:4.2]{Claim 4.2} [The Abel Theorem on the Cauchy product].

\begin{claim}[The Abel Theorem on the Cauchy product]\label{claim:7.9}
Let $\sum_{k=0}^{\infty}a_k$ and $\sum_{k=0}^{\infty}b_k$ be two convergent series with sums $a$ and $b$ respectively.
If their Cauchy product $\sum_{k=0}^{\infty}c_k$, where
\[
c_k=a_0b_k+a_1b_{k-1}+\cdots+a_kb_0,
\qquad \forall k\in\N\cup\{0\},
\]
is a convergent series with the sum $c$, then $c=a\cdot b$.
\end{claim}

\begin{proof}
Let us consider the series:
\[
f(x):=\sum_{k=0}^{\infty}a_kx^k,
\qquad
g(x):=\sum_{k=0}^{\infty}b_kx^k,
\qquad
h(x):=\sum_{k=0}^{\infty}c_kx^k.
\]
Let $R_1$, $R_2$, and $R_3$ be the radii of convergence of $f(x)$, $g(x)$, and $h(x)$ respectively.
Since
\[
f(1)=\sum_{k=0}^{\infty}a_k1^k=a,
\qquad
g(1)=\sum_{k=0}^{\infty}b_k1^k=b,
\qquad
h(1)=\sum_{k=0}^{\infty}c_k1^k=c,
\]
due to \hyperref[def:7.2]{Def 7.2}, we have $R_1\geq1$, $R_2\geq1$, and $R_3\geq1$.
Therefore, due to \hyperref[item:claim:7.3:1]{Item 1)} (if $R_i>1$) or \hyperref[item:claim:7.3:2]{Item 2)} (if $R_i=1$) of \hyperref[claim:7.3]{Claim 7.3}, the series $f(x)$, $g(x)$, and $h(x)$ converge uniformly on the interval $[0,1]$.

Therefore, due to \hyperref[claim:6.11]{Claim 6.11} [Continuity of a Series], the functions $f(x)$, $g(x)$, and $h(x)$ are continuous on $[0,1]$ and we have
\[
\lim_{x\to1-0}f(x)\overset{1}{=}f(1)=a,
\qquad
\lim_{x\to1-0}g(x)\overset{1}{=}g(1)=b,
\qquad
\lim_{x\to1-0}h(x)\overset{1}{=}h(1)=c.
\tag{6}\label{eq:7:6:1}
\]
\hyperref[eq:7:6:1]{Equalities 1)} hold since the functions are continuous at $x=1$.

Since $R_1\geq1$ and $R_2\geq1$, due to \hyperref[item:claim:7.1:1]{Item 1)} of \hyperref[claim:7.1]{Claim 7.1}, the series $f(x)=\sum_{k=0}^{\infty}a_kx^k$ and $g(x)=\sum_{k=0}^{\infty}b_kx^k$ converge absolutely for every $x\in[0,1)$.

Note that, for every $k\in\N\cup\{0\}$, $\forall x\in\R$, we have
\begin{align*}
c_kx^k
&=(a_0x^0)(b_kx^k)+(a_1x^1)(b_{k-1}x^{k-1})+\cdots+(a_kx^k)(b_0x^0).
\end{align*}
That is, due to \hyperref[def:4.1]{Def 4.1},
\[
h(x)=\sum_{k=0}^{\infty}c_kx^k
\]
is the Cauchy product of the series
\[
\sum_{k=0}^{\infty}a_kx^k=f(x)
\qquad\text{and}\qquad
\sum_{k=0}^{\infty}b_kx^k=g(x).
\]
Therefore, due to \hyperref[claim:4.1]{Claim 4.1} [The Mertens Theorem], we have
\[
h(x)=f(x)\cdot g(x),
\qquad \text{for every }x\in[0,1).
\tag{7}\label{eq:7:7:1}
\]
To finish the proof, note that:
\[
c
\overset{1}{=}
\lim_{x\to1-0}h(x)
\overset{2}{=}
\lim_{x\to1-0}\bigl(f(x)g(x)\bigr)
\overset{3}{=}
\left(\lim_{x\to1-0}f(x)\right)
\left(\lim_{x\to1-0}g(x)\right)
\overset{4}{=}
ab.
\]
Here, 1) and 4) are due to \hyperref[eq:7:6:1]{Equalities (6)}, 2) is due to \hyperref[eq:7:7:1]{Equality (7)}, and 3) is due to the \hyperref[eq:7:7:1]{Equality (7)} for every $x<1$ and we can assume that $x>0$.
\end{proof}

\begin{claim}\label{claim:7.10}
Let $\sum_{k=0}^{\infty}a_kx^k$ be a power series with the radius of convergence $R>0$.
Then
\[
\forall p\in(0,R)\quad \exists C>0\quad \forall k\in\N\cup\{0\}
\quad\to\quad
|a_k|\leq\frac{C}{p^k}.
\]
\end{claim}

\begin{proof}
Since $p\in(0,R)$, due to \hyperref[item:claim:7.1:1]{Item 1)} of \hyperref[claim:7.1]{Claim 7.1}, the series
\[
\sum_{k=0}^{\infty}|a_k|p^k
\]
converges. Therefore, due to \hyperref[claim:1.5]{Claim 1.5},
\[
\lim_{k\to\infty}|a_k|p^k=0.
\]
It is well-known (see Calculus-1) that every convergent sequence is bounded. Thus we have
\[
\exists C>0\quad \forall k\in\N\cup\{0\}
\quad\to\quad
|a_k|p^k<C.
\]
Since $p>0$, the claim is proved.
\end{proof}
